% Generated by tools/build.js from content/ -- edit the YAML, not this file.
\documentclass[10pt,letterpaper]{article}
\usepackage[T1]{fontenc}
\usepackage[utf8]{inputenc}
\usepackage{lmodern}
\usepackage{textcomp}
\usepackage[margin=0.6in]{geometry}
\usepackage{enumitem}
\usepackage[hidelinks]{hyperref}
\hypersetup{pdftitle={Jesús M. Rueda-Becerril, Ph.D.}, pdfauthor={Jesús M. Rueda-Becerril}}
% map every glyph (including ligatures) to Unicode so text extraction is exact
\input{glyphtounicode}
\pdfgentounicode=1

\pagestyle{empty}
\setlength{\parindent}{0pt}
\setlength{\parskip}{0pt}
\setlength{\emergencystretch}{3em}
\setlist[itemize]{leftmargin=1.3em, topsep=1pt, itemsep=0.5pt, parsep=0pt}
\newcommand{\cvsection}[1]{%
  \vspace{5pt}{\large\bfseries #1}\par\nopagebreak\vspace{2pt}\hrule
  \nopagebreak\vspace{3pt}\nopagebreak}
\newcommand{\cvsubsection}[1]{%
  \vspace{3pt}{\bfseries #1}\par\nopagebreak\vspace{2pt}\nopagebreak}
\newcommand{\entrygap}{\vspace{3pt}}
% left text, then right text flush right on the same line if there is at
% least 1em to spare, otherwise flush right on the next line (TeXbook \signed)
\newcommand{\leftright}[2]{{#1\unskip\nobreak\hfil\penalty50\hskip1em\hbox{}%
  \nobreak\hfil\mbox{#2}\parfillskip=0pt\par}}

\begin{document}

\begin{center}
  {\LARGE\bfseries Jesús M. Rueda-Becerril, Ph.D.}\\[3pt]
  Seattle, WA \textbar{} \href{mailto:jm.ruebe@gmail.com}{jm.ruebe@gmail.com} \textbar{} +1 (765) 430-2330 \textbar{} \href{https://linkedin.com/in/jeruebe}{linkedin.com/in/jeruebe}\\
  \href{https://altjerue.github.io}{altjerue.github.io} \textbar{} \href{https://github.com/altjerue}{github.com/altjerue} \textbar{} \href{https://orcid.org/0000-0003-1988-1912}{orcid.org/0000-0003-1988-1912} \textbar{} \href{https://scholar.google.com/citations?user=hrld5zgAAAAJ}{Google Scholar}
\end{center}

\cvsection{Professional Summary}
Computational physicist who turns complex physical models into fast, tested software and turns large scientific datasets into calibrated models. Ph.D. in Physics and ten years building simulation codes and the data pipelines around them: author of Paramo, an open-source Fortran/OpenMP radiative-transfer library used by other research groups; led a 7-person upgrade of the PatchworkMHD black-hole simulation code on the Frontera supercomputer; mapped wetland probability across a 2,160 km² watershed with random-forest models. Nine peer-reviewed articles and Co-I on a funded NASA Fermi Guest Investigator grant.\par

\cvsection{Skills}
\textbf{ML \& Data:} Scientific ML, PyTorch, TensorFlow, scikit-learn, NumPy, SciPy, Pandas, Astropy, random forest, SVM, neural networks, gradient-based optimization, parameter estimation, uncertainty quantification, synthetic dataset generation\par
\textbf{Scientific Computing \& HPC:} Numerical PDEs (Fokker–Planck, advection–diffusion), radiative transfer, relativistic MHD, Lagrangian particle dynamics, finite-difference/finite-volume methods, numerical linear algebra, OpenMP, MPI, OpenACC, SLURM, PBS, HDF5\par
\textbf{Geospatial \& Remote Sensing:} Sentinel-1/2 multispectral imagery, DEM terrain analysis, raster/vector pipelines, QGIS, geospatial feature engineering\par
\textbf{Languages:} Python, C/C++, Fortran, Rust, R, SQL, Java, Julia, Shell\par
\textbf{Tools:} Git, Docker, Azure, Jenkins, OpenShift, Kafka, Splunk, Jira, ParaView, VisIt\par

\cvsection{Experience}
\leftright{\textbf{Independent Contractor}}{Jan 2026 – Present}\nopagebreak
\leftright{\textit{Mercor}}{Remote}\nopagebreak
Physics subject-matter expert evaluating frontier language models for a leading AI lab.\par\nopagebreak
\begin{itemize}
  \item{} Reviewed 50+ expert-level problems in physics, hydrodynamics, astrophysics, and cosmology, and wrote prompts that stumped frontier models.
  \item{} Diagnose failure modes in model reasoning and write structured feedback to improve scientific accuracy.
  \item{} Benchmark model performance and quality-check solutions to complex STEM problems.
\end{itemize}
\entrygap

\leftright{\textbf{Spatial Data Scientist}}{May 2025 – Jan 2026}\nopagebreak
\leftright{\textit{TealWaters}}{Seattle, WA}\nopagebreak
Built geospatial ML pipelines and optimized terrain-analysis code for wetland mapping.\par\nopagebreak
\begin{itemize}
  \item{} Mapped wetland probability across the 2,160 km² Skykomish watershed with the Wetland Intrinsic Potential (WIP) random-forest model, reaching >75\% accuracy.
  \item{} Optimized the terrain-modeling code that computes elevation derivatives from 1 m, 3 m, and 10 m DEMs over 20 tiles, improving its computational performance.
  \item{} Engineered features from Sentinel-1/2 multispectral imagery and terrain derivatives for ML land-cover and wetland models; built Python processing tools and QGIS visualizations.
  \item{} Presented results to managers and decision-makers, trained 2 colleagues on the WIP tool, and worked with software engineers to move prototype tools into production.
\end{itemize}
\entrygap

\leftright{\textbf{Independent Researcher \& Open-Source Developer}}{Jan 2024 – May 2025}\nopagebreak
\leftright{\textit{Self-directed}}{Seattle, WA}\nopagebreak
Continued simulation research and open-source development between industry roles.\par\nopagebreak
\begin{itemize}
  \item{} Co-developed Tleco, a Rust/Python toolkit that models radiation from relativistic outflows by solving the Fokker–Planck equation coupled to radiative transfer; published in ApJ (2024) [9].
  \item{} Developing WindsOfChange with the UW Forest Mycobiome Lab: a C++/R library that simulates wind-driven spore dispersal over forested, hilly terrain, parallelized with RcppParallel and OpenMP.
  \item{} Invited talk on machine learning in environmental science at the University of Washington (2024).
\end{itemize}
\entrygap

\leftright{\textbf{Software Engineer}}{Apr 2022 – Jan 2024}\nopagebreak
\leftright{\textit{Paychex}}{Remote}\nopagebreak
Built and tested data-integration services between enterprise systems and Oracle E-Business Suite.\par\nopagebreak
\begin{itemize}
  \item{} Built Java Kafka consumers that moved data between enterprise databases, deployed through Jenkins and OpenShift CI/CD pipelines.
  \item{} Created Splunk dashboards tracking the users shared between enterprise systems and Oracle EBS, with alerts that flagged production incidents.
  \item{} Wrote full test suites for 5+ Python services in the CI pipeline, keeping production defect rates low.
  \item{} Prepared, validated, and analyzed Oracle EBS data with SQL.
\end{itemize}
\entrygap

\leftright{\textbf{Postdoctoral Research Associate}}{Feb 2021 – Apr 2022}\nopagebreak
\leftright{\textit{Rochester Institute of Technology}}{Rochester, NY}\nopagebreak
HPC simulations of supermassive black-hole binaries in an NSF-funded project.\par\nopagebreak
\begin{itemize}
  \item{} Led a 7-person team upgrading PatchworkMHD, a C code for magnetohydrodynamic simulations of supermassive black-hole binaries.
  \item{} Added black-hole spin to the simulations with under 5\% runtime overhead.
  \item{} Designed the simulation experiments and ran and benchmarked the code on TACC's Frontera supercomputer, identifying performance optimization opportunities.
  \item{} Calibrated Paramo model parameters against Fermi-LAT gamma-ray spectra with gradient-based optimization, and supervised a graduate student applying gradient descent to the same problem.
  \item{} Co-authored 2 papers from a multi-institution study of binary neutron-star mergers [6, 7] and a third on blazar heating and cooling [8].
\end{itemize}
\entrygap

\leftright{\textbf{Postdoctoral Research Fellow}}{Oct 2018 – Nov 2020}\nopagebreak
\leftright{\textit{Purdue University}}{West Lafayette, IN}\nopagebreak
Radiative-transfer modeling of relativistic jets (blazars and gamma-ray burst afterglows).\par\nopagebreak
\begin{itemize}
  \item{} Architected Paramo, an open-source Fortran 95 library that solves the Fokker–Planck equation coupled to radiative transfer; parallelized it with OpenMP for a 20× speedup over the serial code on 10 cores.
  \item{} Paramo has been adopted by other research groups, including in Combi \& Siegel (2023).
  \item{} Co-I on two NASA Fermi Guest Investigator proposals; one was funded (\$68K) and supported the finding that blazar jets are launched with similar energy per baryon, independently of their power [5].
  \item{} Added numerical-integration and OpenMP features to Paramo and built Python tools to compute energy losses and spectra of high-energy particles in gamma-ray burst afterglows [4].
  \item{} Mentored 3 graduate students (1 M.S., 2 Ph.D.), co-authoring papers with each.
  \item{} Wrote Python scripts for statistical modeling of the COVID-19 outbreak in Mexico with a multidisciplinary group, and co-wrote Spanish-language infographics and blog posts against misinformation.
\end{itemize}
\entrygap

\leftright{\textbf{Postdoctoral Research Fellow}}{Jan 2018 – Sep 2018}\nopagebreak
\leftright{\textit{Universidad Michoacana de San Nicolás de Hidalgo}}{Morelia, Mexico}\nopagebreak
\begin{itemize}
  \item{} Built a Python pipeline that turned GRTrans black-hole simulation images into training data for an SVM classifier of black-hole spin (grayscale conversion, intensity normalization, spatial feature extraction).
  \item{} Released an open-source Python tool for computing radiative-transfer spectra and light curves in relativistic astrophysics.
  \item{} Organized an HDF5 workshop to train graduate students in high-volume data storage.
\end{itemize}
\entrygap

\leftright{\textbf{Graduate Research Assistant}}{Oct 2011 – Jul 2017}\nopagebreak
\leftright{\textit{Universitat de València}}{Valencia, Spain}\nopagebreak
\begin{itemize}
  \item{} Automated runs of the C-SPEV radiative-transfer code with Shell and Python pipelines: 10–60 simulations per project and tens of GB of HDF5 output, with data-quality checks for downstream analysis.
  \item{} Analyzed multi-wavelength observations from the NASA Fermi-LAT telescope and the Very Long Baseline Array (VLBA), including cross-matching sources across catalogs by sky coordinates and fitting log-parabola spectra.
  \item{} Fit C-SPEV simulation models to telescope data, identifying spectral and light-curve patterns that quantify the magnetization of plasma in blazars.
  \item{} Extended C-SPEV to compute discrete and continuous spectra from particle distributions of arbitrary shape, without affecting simulation runtime.
  \item{} Published 2 first-author papers in MNRAS [2, 3].
\end{itemize}
\entrygap


\cvsection{Education}
\leftright{\textbf{Ph.D. in Physics}}{Oct 2011 – Jul 2017}\nopagebreak
\leftright{\textit{Universitat de València}}{Valencia, Spain}\nopagebreak
Excellent \textit{cum laude}\par\nopagebreak
Advisors: Prof. Miguel A. Aloy and Dr. Petar Mimica\par\nopagebreak
Thesis: \href{http://roderic.uv.es/handle/10550/60003}{Numerical treatment of radiation processes in the internal shocks of magnetized relativistic outflows}\par
\entrygap

\leftright{\textbf{M.Sc. in Physics}}{Aug 2009 – Sep 2011}\nopagebreak
\leftright{\textit{Universidad Michoacana de San Nicolás de Hidalgo}}{Morelia, Mexico}\nopagebreak
Advisor: Prof. José A. Cervera\par\nopagebreak
Thesis: \textit{Study of TOV stars with the SPH method}\par
\entrygap

\leftright{\textbf{B.Sc. in Physics}}{Aug 2004 – Dec 2008}\nopagebreak
\leftright{\textit{Universidad Autónoma del Estado de México}}{Toluca, Mexico}\nopagebreak
Advisor: Prof. Francisco S. Guzmán\par\nopagebreak
Thesis: \textit{Numerical solution of null geodesics for the generation of gravitational lenses produced by spherically-symmetric and static spacetimes}\par\nopagebreak
Lic. Juan Josafat Pichardo Cruz Award\par
\entrygap


\cvsection{Publications}
\cvsubsection{Articles (9 published, 3 first-author; 3 in preparation)}
\begin{itemize}[label={}, leftmargin=1.5em, itemindent=-1.5em, labelwidth=0pt, labelsep=0pt]
  \item{} [12] Winkowski, S., \textbf{Rueda-Becerril, J. M.}, \& Willing, C. \textit{Using Numerical Simulations of Wind on Hilly Terrain to Model Propagule Dispersal}, in prep.
  \item{} [11] Holzman, L., \textbf{Rueda-Becerril, J. M.}, Creamer, C., \& Willing, C. \textit{Random Forest Link Landscape Position to Drought Resilience in Giant Sequoia Forests}, in prep.
  \item{} [10] \textbf{Rueda-Becerril, J. M.}, Winkowski, S., \& Willing, C. \textit{WindsOfChange: A New Toolkit for Modeling Spore Dispersal in Hilly Terrains}, in prep.
  \item{} [9] Davis, Z., \textbf{Rueda-Becerril, J. M.}, \& Giannios, D. \textit{Tleco: A Toolkit for Modeling Radiative Signatures from Relativistic Outflows}, \href{https://doi.org/10.3847/1538-4357/ad8bc2}{ApJ 976, 182 (2024)}, \href{https://arxiv.org/abs/2405.17581}{arXiv:2405.17581}.
  \item{} [8] Davis, Z., \textbf{Rueda-Becerril, J. M.}, \& Giannios, D. \textit{Balancing Turbulent Heating with Radiative Cooling in Blazars}, \href{https://doi.org/10.1093/mnras/stac1282}{MNRAS 513, 5766–5779 (2022)}, \href{https://arxiv.org/abs/2201.07790}{arXiv:2201.07790}.
  \item{} [7] Lopez-Armengol, F. G., Etienne, Z. B., [...], \textbf{Rueda-Becerril, J. M.}, [...] \textit{Handing off the outcome of binary neutron star mergers for accurate and long-term postmerger simulations}, \href{https://doi.org/10.1103/PhysRevD.106.083015}{Phys. Rev. D 106, 083015 (2022)}, \href{https://arxiv.org/abs/2112.09817}{arXiv:2112.09817}.
  \item{} [6] Murguia-Berthier, A., Noble, S., [...], \textbf{Rueda-Becerril, J. M.}, [...] \textit{HARM3D+NUC: A New Method for Simulating the Post-merger Phase of Binary Neutron Star Mergers with GRMHD, Tabulated EOS, and Neutrino Leakage}, \href{https://doi.org/10.3847/1538-4357/ac1119}{ApJ 919, 95 (2021)}, \href{https://arxiv.org/abs/2106.05356}{arXiv:2106.05356}.
  \item{} [5] \textbf{Rueda-Becerril, J. M.}, Harrison, A. O., \& Giannios, D. \textit{Blazar jets launched with similar energy per baryon, independently of their power}, \href{https://doi.org/10.1093/mnras/staa3925}{MNRAS 501, 4092–4102 (2021)}, \href{https://arxiv.org/abs/2009.02273}{arXiv:2009.02273}.
  \item{} [4] Zhang, H., Christie, I., Petropoulou, M., \textbf{Rueda-Becerril, J. M.}, \& Giannios, D. \textit{Inverse Compton Signatures of Gamma-Ray Burst Afterglows}, \href{https://doi.org/10.1093/mnras/staa1583}{MNRAS 496, 974–986 (2020)}, \href{https://arxiv.org/abs/1910.14049}{arXiv:1910.14049}.
  \item{} [3] \textbf{Rueda-Becerril, J. M.}, Mimica, P., \& Aloy, M. A. \textit{On the influence of a hybrid thermal–non-thermal distribution in the internal shocks model for blazars}, \href{https://doi.org/10.1093/mnras/stx476}{MNRAS 468, 1169–1182 (2017)}, \href{https://arxiv.org/abs/1612.06383}{arXiv:1612.06383}.
  \item{} [2] \textbf{Rueda-Becerril, J. M.}, Mimica, P., \& Aloy, M. A. \textit{The influence of the magnetic field on the spectral properties of blazars}, \href{https://doi.org/10.1093/mnras/stt2335}{MNRAS 438, 1856–1869 (2014)}, \href{https://arxiv.org/abs/1310.5441}{arXiv:1310.5441}.
  \item{} [1] Guzmán, F. S., \& \textbf{Rueda-Becerril, J. M.} \textit{Spherical boson stars as black hole mimickers}, \href{https://doi.org/10.1103/PhysRevD.80.084023}{Phys. Rev. D 80, 084023 (2009)}, \href{https://arxiv.org/abs/1009.1250}{arXiv:1009.1250}.
\end{itemize}
\cvsubsection{Conference Proceedings}
\begin{itemize}[label={}, leftmargin=1.5em, itemindent=-1.5em, labelwidth=0pt, labelsep=0pt]
  \item{} [P5] \textbf{Rueda-Becerril, J. M.} \textit{A numerical approach for radiative cooling in relativistic outflows}, \href{https://doi.org/10.1002/asna.202113919}{Astron. Nachr. 342, 277–282 (2021)}, 9th International Workshop on Astronomy and Relativistic Astrophysics, \href{https://arxiv.org/abs/2011.13797}{arXiv:2011.13797}.
  \item{} [P4] \textbf{Rueda-Becerril, J. M.}, Harrison, A. O., \& Giannios, D. \textit{The blazar sequence revised}, \href{https://doi.org/10.1002/asna.202113895}{Astron. Nachr. 342, 147–152 (2021)}, 9th International Workshop on Astronomy and Relativistic Astrophysics, \href{https://arxiv.org/abs/2011.13805}{arXiv:2011.13805}.
  \item{} [P3] \textbf{Rueda-Becerril, J. M.}, Mimica, P., \& Aloy, M. A. \textit{Numerical simulations of the internal shock model in magnetized relativistic jets of blazars}, \href{https://doi.org/10.22323/1.233.0159}{PoS(SWIFT 10) 233, 159 (2014)}, \href{https://arxiv.org/abs/1502.07882}{arXiv:1502.07882}.
  \item{} [P2] \textbf{Rueda-Becerril, J. M.}, Mimica, P., Aloy, M. A., \& Aloy, C. \textit{Numerical study of broadband spectra caused by internal shocks in magnetized relativistic jets of blazars}, \href{https://doi.org/10.1051/epjconf/20136102007}{EPJ Web Conf. 61, 02007 (2013)}, \href{https://arxiv.org/abs/1309.4612}{arXiv:1309.4612}.
  \item{} [P1] Mimica, P., Aloy, M. A., \textbf{Rueda-Becerril, J. M.}, Tabik, S., \& Aloy, C. \textit{Numerical simulations of dynamics and emission from relativistic astrophysical jets}, \href{https://doi.org/10.1088/1742-6596/454/1/012001}{J. Phys.: Conf. Ser. 454, 012001 (2013)}, \href{https://arxiv.org/abs/1211.1794}{arXiv:1211.1794}.
\end{itemize}

\cvsection{Grants and Fellowships}
\cvsubsection{Research Grants}
\leftright{\textbf{NASA Fermi Guest Investigator Program, Cycle 12}}{2019}\nopagebreak
\leftright{\textit{A simple model to understand the blazar sequence}}{Grant \#121077, \$68K}\nopagebreak
PI: D. Giannios; Co-I: J. M. Rueda-Becerril\par
\entrygap

\cvsubsection{Fellowships}
\leftright{\textbf{International Postdoctoral Fellowship}}{Oct 2018 – Nov 2020}\nopagebreak
\textit{National Council of Science and Technology (CONACyT), Mexico}\par
\entrygap

\leftright{\textbf{Postdoctoral Fellowship}}{Jan 2018 – Sep 2018}\nopagebreak
\textit{Secretariat of Public Education, Mexico (Program for the Professional Development of Higher Education Institutions)}\par
\entrygap

\leftright{\textbf{Study-Abroad Fellowship}}{Sep 2014 – Aug 2016}\nopagebreak
\textit{National Council of Science and Technology (CONACyT), Mexico}\par
\entrygap

\leftright{\textbf{Santiago Grisolía Fellowship}}{Oct 2011 – Jun 2014}\nopagebreak
\textit{Council of Education, Research, Culture and Sport, Valencian Community, Spain}\par
\entrygap

\leftright{\textbf{M.Sc. Fellowship}}{Sep 2009 – Aug 2011}\nopagebreak
\textit{CONACyT, at the Institute of Physics and Mathematics, Universidad Michoacana de San Nicolás de Hidalgo}\par
\entrygap

\leftright{\textbf{Summer Research Fellowship (XVII Summer of Scientific Research)}}{Jun 2007 – Aug 2007}\nopagebreak
\textit{Mexican Academy of Sciences}\par
\entrygap


\cvsection{Software and Projects}
\leftright{\textbf{Blazar Population ML}}{Jan 2026 – Present}\nopagebreak
\leftright{\textit{Scientific ML (Python, scikit-learn, PyTorch)}}{\href{https://github.com/altjerue/ilwikak-ml}{github.com/altjerue/ilwikak-ml}}\nopagebreak
\begin{itemize}
  \item{} Applying classical ML and Bayesian inference to the Fermi-LAT 4LAC-DR3 catalog (3,407 sources) to recover physical jet parameters from multi-wavelength data.
  \item{} Building a neural network that maps physical jet parameters to observable spectral energy distributions accurately enough to compare with Tleco results.
  \item{} Forecasting blazar gamma-ray flares with a time-series foundation model fine-tuned on Tleco-generated synthetic light curves.
\end{itemize}
\entrygap

\leftright{\textbf{WindsOfChange}}{Aug 2024 – Present}\nopagebreak
\leftright{\textit{Lagrangian particle-dispersion library (C++/R)}}{\href{https://github.com/UW-ForestMycobiomeLab/WindSporesHills}{github.com/UW-ForestMycobiomeLab/WindSporesHills}}\nopagebreak
\begin{itemize}
  \item{} Simulates particle dispersion, atmospheric turbulence, and flow over hilly terrain with forest canopy, based on analytical models (Massman \& Weil 1999; Finnigan \& Belcher 2004); designed to scale to tens of thousands of particles over kilometer-scale domains.
\end{itemize}
\entrygap

\leftright{\textbf{Tleco}}{Jan 2024 – Sep 2024}\nopagebreak
\leftright{\textit{Radiative-transfer toolkit (Rust/Python)}}{\href{https://github.com/zkdavis/Tleco}{github.com/zkdavis/Tleco}}\nopagebreak
\begin{itemize}
  \item{} Numerical PDE solvers, particle-distribution algorithms, and spectral calculations for relativistic plasmas; successor to Paramo [9].
\end{itemize}
\entrygap

\leftright{\textbf{Paramo}}{Oct 2018 – Apr 2022}\nopagebreak
\leftright{\textit{Radiative-transfer library (Fortran 95, OpenMP)}}{\href{https://github.com/altjerue/paramo}{github.com/altjerue/paramo}}\nopagebreak
\begin{itemize}
  \item{} Fokker–Planck and radiative-transfer solver for relativistic particle distributions; 20× OpenMP speedup on 10 cores; used by other groups, including Combi \& Siegel (2023).
\end{itemize}
\entrygap

\leftright{\textbf{PatchworkMHD}}{Feb 2021 – Apr 2022}\nopagebreak
\textit{Large-scale MHD simulation code (NSF project)}\par\nopagebreak
\begin{itemize}
  \item{} Led a 7-person upgrade of this production C code for HPC environments, adding black-hole spin with under 5\% runtime overhead.
\end{itemize}
\entrygap


\cvsection{Awards}
\leftright{\textbf{Marcos Moshinsky Award for Best Poster}}{Sep 2020}\nopagebreak
\textit{9th International Workshop on Astronomy and Relativistic Astrophysics (IWARA 2020), video conference}\par
\entrygap

\leftright{\textbf{Lic. Juan Josafat Pichardo Cruz Award}}{2009}\nopagebreak
\textit{Universidad Autónoma del Estado de México}\par\nopagebreak
For defending a licentiate thesis within a year of completing the undergraduate coursework.\par
\entrygap


\cvsection{Presentations}
\cvsubsection{Invited Talks}
\leftright{\textbf{The Role of Machine Learning in Environmental Science}}{Mar 2024}\nopagebreak
\leftright{\textit{Seminar on Climate Change Ecology, University of Washington}}{Seattle, WA}
\entrygap

\leftright{\textbf{Morphology of the spectra from numerical simulations of the internal shocks model for blazars}}{Feb 2019}\nopagebreak
\leftright{\textit{Astrophysics Seminar, Purdue University}}{West Lafayette, IN}
\entrygap

\leftright{\textbf{Numerical simulations of the internal shocks model in magnetized relativistic jets of blazars}}{Jun 2018}\nopagebreak
\leftright{\textit{DATA Group Weekly Seminar, Instituto de Astronomía, UNAM}}{Mexico City, Mexico}
\entrygap

\leftright{\textbf{Numerical treatment of non-thermal radiation in the internal shocks model for blazars}}{Mar 2018}\nopagebreak
\leftright{\textit{Weekly Seminar, Instituto de Física y Matemáticas, Universidad Michoacana de San Nicolás de Hidalgo}}{Morelia, Mexico}
\entrygap

\leftright{\textbf{Numerical simulations of the internal shock model in magnetized relativistic jets of blazars}}{Oct 2014}\nopagebreak
\leftright{\textit{IVICFA's Fridays: Computation in Physics, IFIC}}{Paterna, Spain}
\entrygap

\cvsubsection{Contributed Talks}
\leftright{\textbf{Simulations of supermassive binary black holes accretion dynamics in the spinning case}}{Apr 2022}\nopagebreak
\leftright{\textit{APS April Meeting 2022}}{New York, NY}
\entrygap

\leftright{\textbf{A numerical approach to the Klein-Nishina corrections of radiative cooling in relativistic outflows}}{Apr 2021}\nopagebreak
\leftright{\textit{APS April Meeting 2021}}{Online}
\entrygap

\leftright{\textbf{The blazar sequence revised}}{Sep 2020}\nopagebreak
\leftright{\textit{9th International Workshop on Astronomy and Relativistic Astrophysics}}{Online}\nopagebreak
Video: \href{https://www.youtube.com/watch?v=BAZNWLNT69M}{youtube.com/watch?v=BAZNWLNT69M}\par
\entrygap

\leftright{\textbf{Influence of the magnetic field on the spectral properties of blazars in the internal shocks scenario}}{Jun 2014}\nopagebreak
\leftright{\textit{Extreme-Astrophysics in an Ever-Changing Universe}}{Ierápetra, Greece}
\entrygap

\leftright{\textbf{Numerical study of broadband spectra caused by internal shocks in magnetized relativistic jets}}{Jul 2013}\nopagebreak
\leftright{\textit{XXXIV Biennial Meeting of the Royal Spanish Society of Physics}}{Valencia, Spain}
\entrygap

\cvsubsection{Posters}
\leftright{\textbf{A numerical approach for radiative cooling in relativistic outflows}}{Sep 2020}\nopagebreak
\leftright{\textit{9th International Workshop on Astronomy and Relativistic Astrophysics}}{Online}\nopagebreak
Marcos Moshinsky Award for Best Poster. Video: \href{https://www.youtube.com/watch?v=0TJiKg7kOPI}{youtube.com/watch?v=0TJiKg7kOPI}\par
\entrygap

\leftright{\textbf{Numerical simulations of the internal shock model in magnetized relativistic jets of blazars}}{Dec 2014}\nopagebreak
\leftright{\textit{Swift: 10 Years of Discovery}}{Rome, Italy}
\entrygap

\leftright{\textbf{Numerical study of broadband spectra caused by internal shocks in magnetized relativistic jets}}{Jun 2013}\nopagebreak
\leftright{\textit{The Innermost Regions of Relativistic Jets and Their Magnetic Fields}}{Granada, Spain}
\entrygap


\cvsection{Teaching and Mentoring}
\leftright{\textbf{Zachary Davis, graduate student}}{2018 – 2022}\nopagebreak
\leftright{\textit{Department of Physics and Astronomy, Purdue University}}{Mentoring}\nopagebreak
Co-authored papers [8, 9]\par
\entrygap

\leftright{\textbf{Amanda O. Harrison, graduate student}}{2018 – 2020}\nopagebreak
\leftright{\textit{Department of Physics and Astronomy, Purdue University}}{Mentoring}\nopagebreak
Co-authored paper [5]\par
\entrygap

\leftright{\textbf{Hao Zhang, graduate student}}{2018 – 2019}\nopagebreak
\leftright{\textit{Department of Physics and Astronomy, Purdue University}}{Mentoring}\nopagebreak
Co-authored paper [4]\par
\entrygap

\leftright{\textbf{Thermodynamics (graduate level)}}{Jun 2018}\nopagebreak
\leftright{\textit{Instituto de Física y Matemáticas, Universidad Michoacana de San Nicolás de Hidalgo}}{Substitute Instructor}
\entrygap


\cvsection{Professional Development}
\leftright{\textbf{High Performance Computing on Frontera}}{May – Jun 2021}\nopagebreak
\leftright{\textit{Jason Allison et al., TACC, Austin, TX}}{Lecture series}
\entrygap

\leftright{\textbf{Writing Winning Grants}}{Nov 2019}\nopagebreak
\leftright{\textit{Dr. Lauren Broyles, Purdue University}}{Lecture}
\entrygap

\leftright{\textbf{XSEDE HPC Workshop: Summer Boot Camp}}{Jun 2019}\nopagebreak
\leftright{\textit{John Urbanic, Purdue University}}{Workshop}
\entrygap

\leftright{\textbf{Data Analysis and Machine Learning with Python}}{Feb 2017}\nopagebreak
\leftright{\textit{Dr. Alejandro Torres, Universitat de València}}{Workshop}
\entrygap

\leftright{\textbf{Numerical Relativity Simulations of BBH Coalescence Using the Einstein Toolkit}}{Jul 2016}\nopagebreak
\leftright{\textit{Dr. Vassilios Mewes, Universitat de València}}{Workshop, 8 hours}
\entrygap

\leftright{\textbf{The Universe in the Light of PLANCK and BICEP2}}{May 2014}\nopagebreak
\leftright{\textit{Prof. Nick Mavromatos, Universitat de València}}{Lecture series, 2 credits}
\entrygap

\leftright{\textbf{Dark Matter}}{Sep 2013}\nopagebreak
\leftright{\textit{Prof. Alejandro Ibarra, Universitat de València}}{Lecture series, 2 credits}
\entrygap

\leftright{\textbf{International Cargèse School on Cosmic Accelerators}}{Apr – May 2013}\nopagebreak
\leftright{\textit{Institut d'Études Scientifiques de Cargèse, France}}{Summer school}
\entrygap

\leftright{\textbf{Introduction to C++ Programming}}{Apr 2012}\nopagebreak
\leftright{\textit{Dr. Jacek Generowicz, Universitat de València}}{Workshop, 6 credits}
\entrygap

\leftright{\textbf{Numerical Relativistic Astrophysics}}{Mar – Apr 2012}\nopagebreak
\leftright{\textit{Prof. Luciano Rezzolla, Universitat de València}}{Lecture series, 9 hours}
\entrygap

\leftright{\textbf{Fortran for Scientific Computing}}{Mar 2012}\nopagebreak
\leftright{\textit{HLRS, University of Stuttgart, Germany}}{Workshop, 33 hours}
\entrygap


\cvsection{Certifications}
\leftright{\textbf{Gen AI Intensive}}{\href{https://www.kaggle.com/certification/badges/jmrb103/96}{Kaggle}}\nopagebreak
\textit{Generative AI}\par
\entrygap

\leftright{\textbf{Data Scientist Professional with Python}}{DataCamp}\nopagebreak
\textit{Python programming, data science, data communication, machine learning}\par
\entrygap

\leftright{\textbf{Machine Learning Scientist with Python}}{DataCamp}\nopagebreak
\textit{Machine learning, NLP, deep learning, image processing, big data}\par
\entrygap

\leftright{\textbf{Mathematical Foundations of Machine Learning}}{\href{https://www.udemy.com/certificate/UC-605df108-ae80-4297-8c8f-6bc15b967511/}{Udemy}}\nopagebreak
\textit{Python programming, data science, statistics, machine learning}\par
\entrygap

\leftright{\textbf{Python for Statistical Analysis}}{\href{https://www.udemy.com/certificate/UC-e8557ac8-13f9-41bf-ab46-f196a041b725/}{Udemy}}\nopagebreak
\textit{Python programming, data science, data communication, machine learning}\par
\entrygap


\cvsection{Outreach}
\leftright{\textbf{Los más rápidos y los más furiosos (The Fastest and the Most Furious)}}{Sep 2020}\nopagebreak
\leftright{\textit{Community of Undergraduate Physics Students, Juárez Autonomous University of Tabasco}}{Online talk}
\entrygap

\leftright{\textbf{Una simulación de la física y la astrofísica (A Simulation of Physics and Astrophysics)}}{Aug 2020}\nopagebreak
\leftright{\textit{Community of Undergraduate Physics Students, Juárez Autonomous University of Tabasco}}{Online talk}
\entrygap

\leftright{\textbf{ANITA y la teoría de los universos paralelos (ANITA and the Theory of Parallel Universes)}}{May 2020}\nopagebreak
\leftright{\textit{Científicos Mexicanos en el Extranjero}}{Blog post}\nopagebreak
\href{https://mexiciencia.github.io/post/anita/}{mexiciencia.github.io/post/anita}\par
\entrygap

\leftright{\textbf{¿Qué es el modelo SIR? (What is the SIR Model?)}}{May 2020}\nopagebreak
\leftright{\textit{Científicos Mexicanos en el Extranjero}}{Blog post}\nopagebreak
\href{https://mexiciencia.github.io/post/modelo-sir/}{mexiciencia.github.io/post/modelo-sir}\par
\entrygap

\leftright{\textbf{Evolución del brote epidémico de COVID-19 (Evolution of the COVID-19 Epidemic Outbreak)}}{Apr 2020}\nopagebreak
\leftright{\textit{Científicos Mexicanos en el Extranjero}}{Blog post}\nopagebreak
Data analysis and modeling collaborator. \href{https://mexiciencia.github.io/post/covid19/}{mexiciencia.github.io/post/covid19}\par
\entrygap

\leftright{\textbf{Annual Department of Physics and Astronomy Poster Event}}{Nov 2019}\nopagebreak
\leftright{\textit{Purdue University}}{3 posters}
\entrygap

\leftright{\textbf{Post-Doc Panel Q\&A: What Happens When We Complete Our PhDs?}}{Apr 2019}\nopagebreak
\leftright{\textit{Department of Physics and Astronomy, Purdue University}}{Panelist}
\entrygap

\leftright{\textbf{Annual Department of Physics and Astronomy Poster Event}}{Nov 2018}\nopagebreak
\leftright{\textit{Purdue University}}{Poster}
\entrygap

\leftright{\textbf{¿Decía Einstein la verdad? (Was Einstein Telling the Truth?)}}{Mar 2009}\nopagebreak
\leftright{\textit{Facultad de Ciencias, Universidad Autónoma del Estado de México}}{Talk}
\entrygap


\cvsection{Service}
\leftright{\textbf{Organizing contributor}}{Dec 2012}\nopagebreak
\leftright{\textit{X Scientific Meeting of the Spanish Astronomical Society}}{Valencia, Spain}
\entrygap

\leftright{\textbf{Student representative, Governing Council of the Faculty of Sciences}}{Aug 2007 – May 2009}\nopagebreak
\leftright{\textit{Universidad Autónoma del Estado de México}}{Toluca, Mexico}
\entrygap


\cvsection{Languages}
Spanish (native), English (full professional proficiency), Catalan (basic), Nahuatl (basic)\par

\end{document}
